\documentclass[book.tex]{subfiles}

\begin{document}

\pagestyle{empty}
\vspace*{36 pt}
\begin{center}\Huge\bfseries\itshape
Part 4: Quantum Differentiable Programming
\end{center}
%\newpage

We will now study the use of quantum computing techniques and quantum differentiable programming as a tool for understanding physical systems.

One of the greatest scientific discoveries of the last few decades was the realization that quantum mechanical systems could perform interesting computations for which there exists no  efficient classical analogue. 

At a very high level, the idea is that quantum entanglement provides a new computational resource which has no classical equivalent. By exploiting this resource carefully, quantum algorithms are in some cases able to perform dramatically more efficient calculations than the best known classical equivalent.

For decades, quantum computers remained something of an intellectual curiosity, but this state of affairs has changed drastically in the last few years with the advent of "quantum supremacy." Briefly put, quantum computers are now capable of performing (some not very interesting) calculations faster than any existing classical computer. 

Throughout this book, we have used the new technology of differentiable programs, so you might reasonably wonder whether quantum computing will one day make differentiable programming obsolete. Of course only the future can answer this question completely, but we strongly suspect that differentiable programs and quantum computations will play complementary and mutually beneficial roles in the years to come.

%\subfile{chapters/part4/bloch_sphere}
%\subfile{chapters/part4/quantum_computing_and_quantum_circuits}
%\subfile{chapters/part4/variational_quantum_algorithms}
%\subfile{chapters/part4/quantum_programming}
%\subfile{chapters/part4/quantum_fourier_transform}
%\subfile{chapters/part4/shors_algorithm}
%\subfile{chapters/part4/quantum_machine_learning}

\end{document}